\section{机械部分版本迭代过程记录}
\begin{longtable}[htbp!]{ m{2cm}<{\centering}  m{9.5cm}<{\centering}m{3cm}<{\centering} }
\hline

版本号  &   功能详细说明    &   完成时间    \\
\hline
\endfirsthead
\hline

版本号  &   功能详细说明    &   完成时间    \\
\hline
\endhead
\hline
\multicolumn{2}{r}{接下页}
\endfoot
\caption{机械部分版本迭代过程记录}
\label{tab:机械部分版本迭代过程记录}
\endlastfoot
%   填这里

\hline
 V1.0 & 第1代底盘设计：第一代底盘在22赛季全向轮步兵底盘基础上，仅增大了外部弹仓的容积以追求更大的容弹量。同时增大原本外部弹仓底板的倾斜角度，增加落弹率。   & 2022.11.18 \\

 \hline
 v1.5 & 第1代云台设计：第一代云台在22赛季全向轮步兵云台基础上，增加OAK-D-Pro智能深度相机、灯条、UWB定位模块三位一体固定座。 &2022.11.20 \\

 \hline
 v1.8 & 第1代枪管设计：由于TPU方案长期效果不理想，我们换用双轴承弹丸定心方案，重复度问题得到解决，但双发问题无法解决。 & 2023.01.12  \\

\hline
 V2.0 & 第2代底盘设计：扩大了云台安装座固定铝管的间距与保护框的大小，使得弹仓可修改的空间更大。扩大基座铝管的间距增大了哨兵的容弹量，同时减少了弹丸对底盘重心的影响。优化了中心供弹分层片设计，增大中心供弹齿轮传动的减速比。 & 2023.03.02  \\
 
\hline
v2.5 & 第1.5代云台设计：将OAK-D-Pro智能深度相机更换为MID-360激光雷达，采用雷达里程计以达到更好的避障效果。在1代云台的基础上，新增双枪轮转机构,优化了弹链设计。 & 2023.03.16 \\

\hline
v2.8 & 第2代枪管设计：由于哨兵弹链为活动弹链，弹丸实际轨迹随Pitch轴转动长度会发生细微变化，更换拉簧-打印件定心方案以达到基本解决双发问题的效果。 & 2023.03.20 \\

\hline
v3.0 & 第2代云台设计：更换MID-360激光雷达的固定方式，由原先的在三位一体固定座顶部倒置更换为在三位一体固定座侧方中部侧置，以达到上坡的需求。在第一代双枪轮转机构的基础上，将测速固定板的材质从玻纤更换为6061铝，解决刚性不足问题。& 2023.03.24 \\

\hline

\end{longtable}

\section{电路部分版本迭代过程记录}
\begin{longtable}[htbp!]{ m{2cm}<{\centering}  m{9.5cm}<{\centering}m{3cm}<{\centering} }
\hline

版本号  &   功能详细说明    &   完成时间    \\
\hline
\endfirsthead
\hline

版本号  &   功能详细说明    &   完成时间    \\
\hline
\endhead
\hline
\multicolumn{2}{r}{接下页}
\endfoot
\caption{电路部分版本迭代过程记录}
\label{tab:电路部分版本迭代过程记录}
\endlastfoot
%   填这里

\hline
v1.0 & 第一代高侧开关设计：使用BTS443P芯片作为高侧开关芯片并使用LM5050作为理想二极管防止电流反向 & 2022.10.8 \\ 
\hline
v1.5 & 改进版高侧开关设计：使用LM5060芯片作为高侧控制器芯片理想二极管方案仍采用LM5050芯片 & 2023.1.12\\
\hline
v2.0 & 第二代高侧开关设计：使用LM5060芯片作为高侧控制器并且使用背靠背的NMOS防止电流反向。& 2023.3.12\\
\hline
\end{longtable}

\section{视觉部分版本迭代过程记录}

\begin{longtable}[htbp!]{ m{2cm}<{\centering}  m{9.5cm}<{\centering}m{3cm}<{\centering} }
\hline

版本号  &   功能详细说明    &   完成时间    \\
\hline
\endfirsthead
\hline

版本号  &   功能详细说明    &   完成时间    \\
\hline
\endhead
\hline
\multicolumn{2}{r}{接下页}
\endfoot
\caption{视觉部分版本迭代过程记录}
\label{tab:视觉部分版本迭代过程记录}
\endlastfoot
%   填这里

\hline
v1.0 & 整体流程设计：在视觉标签与相机之间不发生明显位姿变化的条件下（二者的运动处于平行的两条线上）进行了色彩分割、目标提取、内容辨别并分类和位姿估计的设计。& 2023.3.17\\ 
\hline
v2.0 & 在v1.0的基础上，考虑三姿态轴发生变换时，会对内容辨别效果产生影响，采取仿射变换来抗拒姿态变化导致提取的标签的内容的变化，增强分类算法的鲁棒性。& 2023.3.19\\
\hline
\end{longtable}


\section{算法部分版本迭代过程记录}
\begin{longtable}[htbp!]{ m{2cm}<{\centering}  m{9.5cm}<{\centering}m{3cm}<{\centering} }
\hline

版本号  &   功能详细说明    &   完成时间    \\
\hline
\endfirsthead
\hline

版本号  &   功能详细说明    &   完成时间    \\
\hline
\endhead
\hline
\multicolumn{2}{r}{接下页}
\endfoot
\caption{算法版本迭代过程记录}
\label{tab:算法版本迭代过程记录}
\endlastfoot
%   填这里

\hline
v1.0 &  第1代路径规划：第一代路径规划采用OAK-D-Pro深度相机生成点云数据，通过octomap生成稠密点云信息，并将该点云信息投影到二维平面生成代价地图，用于哨兵的局部路径规划；能基本实现自动巡航与避障的功能。& 2022.12.25\\ 
\hline
v1.5 & 第1代哨兵里程计：第一代哨兵里程计采用OAK-D-Pro深度相机自带的视觉里程计，在室内环境下测试下效果尚可。

第1代定位修复：为了解决哨兵里程计在复杂环境与较长时间自动巡航下里程计信息容易发生偏移的问题，我们通用OAK-D-Pro深度相机识别场地上固定位置的标签的位姿，并借此修复自身定位信息。& 2023.02.02\\
\hline

v1.8 & 
第1代状态机：实现了第一版哨兵的状态机，能在手动状态（调试用）与自动状态间通过遥控信号切换；其中自动状态下能实现定点巡逻 & 2023.3.01\\
\hline

v2.0 & 
第2代路径规划：由于深度相机点云范围窄，噪声点多，因此我们将哨兵用于感知的传感器更换为MID-360；同时加入pcl滤波，使生成的代价地图更加清晰可靠，自动避障效果有明显提升  & 2023.03.03\\
\hline

v2.5 & 
第2代哨兵里程计：用MID-360的fast lio定位算法替换了深度相机的视觉里程计，效果有较大提升，且在空旷环境下定位效果同样理想。同时为使雷达兼顾定位算法与代价地图生成算法（这两种算法下所需雷达的消息类型不同），我们实现了零失真的雷达消息类型转换 & 2023.3.10\\
\hline

v3.0 & 
第2代状态机：在原先手动模式与自动模式下，结合自瞄算法，添加了进攻模式的表现效果与触发逻辑，能有效对视野内的敌对单位实施打击

第2代定位修复：在原先视觉标签定位与修复的基础上，加入标签识别数据的过滤，有效减缓了视觉识别下标签抖动与遮挡带来的位置错误修复问题；同时在原先的代码上实现了与apriltag识别的复用

第1代双枪切换：根据比赛要求，实现了双枪的搭载与枪管热量临界时自动切换& 2023.3.13\\
\hline

v3.5 & 
第3代状态机：在受击时能自动进入防御状态，底盘快速自旋，有效减少被打中的概率

第1代爬坡算法：在仿真中通过octomap的高度信息实现斜坡点云数据的过滤，从而在技术上实现哨兵爬坡；同时准备测试Elevation Map在对斜坡点云的过滤效果& 2023.3.29\\
\hline
\end{longtable}